\documentclass[a4paper,11pt]{article}
\usepackage{exam-style-341}
\examinehead{341 农业综合知识三 · 网络技术分章练习}{第05章\quad 局域网与广域网}
\begin{document}

\unitcover{341 农业综合知识三 · 网络技术分章练习}{第05章\quad 局域网与广域网}{练习（共 100 分）}{约 60 分钟}{局域网概述与 IEEE 802 标准；传统以太网与 CSMA/CD；争用期 51.2 $\mu$s 与最短帧长 64 B；截断二进制指数退避；MAC 地址与以太网帧格式；交换机自学习与 VLAN；高速以太网；无线局域网 802.11 与 CSMA/CA；广域网与 PPP、HDLC 等。}

\begin{knowmap}
\textbf{局域网标准：}IEEE 802 系列——802.3 以太网、802.11 无线局域网、802.15 无线个人区域网、802.16 城域网\\[2pt]
\textbf{以太网与 CSMA/CD：}多点接入 · 载波监听 · 碰撞检测；先听后发、边发边听、冲突停止、随机重发\\[2pt]
\textbf{两个关键参数：}10 Mb/s 以太网争用期 $51.2\ \mu$s（$=512$ 比特时间），最短帧长 64 B（$=512$ bit）；帧太短则视为冲突碎片\\[2pt]
\textbf{退避算法：}截断二进制指数退避——第 $i$ 次冲突后取 $k=\min(i,10)$，从 $\{0,1,\dots,2^k-1\}$ 中随机取 $r$，退避 $r$ 倍争用期\\[2pt]
\textbf{MAC 与交换机：}MAC 地址 48 bit；交换机依据帧的源 MAC 地址自学习建立地址表，按目的 MAC 地址转发或过滤，可隔离碰撞域\\[2pt]
\textbf{扩展与高速：}VLAN（802.1Q）划分广播域；100BASE-T、吉比特以太网、10 吉比特以太网；全双工方式下不使用 CSMA/CD\\[2pt]
\textbf{无线与广域网：}802.11 使用 CSMA/CA（可辅以 RTS/CTS 解决隐藏站问题）；广域网常用 PPP、HDLC 等点对点协议
\end{knowmap}

\begin{keyterms}
\textbf{局域网（LAN）}：覆盖范围较小、数据传输速率高、误码率低的计算机网络，通常为一个单位所有。\\
\textbf{以太网}：IEEE 802.3 标准实现的局域网，采用 CSMA/CD 协议和广播信道。\\
\textbf{CSMA/CD}：载波监听多点接入／碰撞检测，以太网使用的媒体访问控制协议。\\
\textbf{争用期}：以太网端到端往返传播时延的 2 倍，10 Mb/s 以太网取 51.2 $\mu$s，又称碰撞窗口。\\
\textbf{最短帧长}：以太网规定的最小有效帧长度为 64 B，小于 64 B 的帧是因冲突而异常中止的无效帧。\\
\textbf{截断二进制指数退避}：发生冲突后按冲突次数从有限范围内随机取数、延迟若干倍争用期再重发的算法。\\
\textbf{MAC 地址}：固化在网络适配器中的 48 bit 硬件地址，用于标识局域网上的站点。\\
\textbf{交换机自学习}：交换机根据收到帧的源 MAC 地址和端口号建立／更新 MAC 地址表的过程。\\
\textbf{VLAN}：虚拟局域网，通过划分逻辑工作组把一个物理局域网划分为多个广播域。\\
\textbf{高速以太网}：速率达到或超过 100 Mb/s 的以太网，如 100BASE-T、吉比特以太网、10 吉比特以太网。\\
\textbf{802.11}：IEEE 无线局域网标准，其 MAC 层使用 CSMA/CA 协议。\\
\textbf{CSMA/CA}：载波监听多点接入／碰撞避免，无线局域网使用的媒体访问控制协议。\\
\textbf{广域网（WAN）}：覆盖范围大、通常借用公用通信网传输的计算机网络。\\
\textbf{PPP}：点对点协议，广泛用于广域网中路由器与路由器、用户与 ISP 之间的点对点链路。
\end{keyterms}

\begin{contrast}{CSMA/CD / CSMA/CA（每行依次为：适用网络·冲突处理·典型标准）}
\citem{CSMA/CD}{有线以太网（总线型、半双工）}{边发送边检测冲突，检测到冲突即停止并退避重发}{IEEE 802.3}
\citem{CSMA/CA}{无线局域网}{无法边发边检测，采用预约信道、确认与退避机制避免冲突}{IEEE 802.11}
\end{contrast}

\begin{contrast}{集线器 / 交换机（每行依次为：工作层次·带宽与碰撞域·工作方式）}
\citem{集线器}{物理层}{所有端口共享带宽，同属一个碰撞域}{半双工，采用 CSMA/CD}
\citem{交换机}{数据链路层}{每个端口独享带宽，可隔离碰撞域}{可全双工，不使用 CSMA/CD}
\end{contrast}

\begin{contrast}{局域网 / 广域网（每行依次为：覆盖范围·核心技术·典型协议）}
\citem{局域网}{小范围（一个单位、一栋楼、一个园区）}{以太网、无线局域网等广播式或交换式技术}{IEEE 802.3、802.11}
\citem{广域网}{大范围（城市、国家、全球）}{点对点链路、分组交换、租用公用通信网}{PPP、HDLC、帧中继、ATM}
\end{contrast}

\begin{contrast}{VLAN / 普通局域网（每行依次为：广播域·隔离能力·配置方式）}
\citem{普通局域网}{一个物理局域网通常是一个广播域}{冲突域可由交换机隔离，广播域不能}{物理连接决定}
\citem{VLAN}{可把一个物理局域网划分为多个广播域}{可隔离广播域、提高安全性和管理灵活性}{通过交换机软件配置划分}
\end{contrast}

\qtypename{一、单项选择题}{10 题，每题 3 分，共 30 分}
\q IEEE 802.3 标准对应的网络是（\quad）
\choicesline{A. 以太网}{B. 无线局域网}{C. 令牌环网}{D. 城域网}

\q CSMA/CD 协议的三个要点是（\quad）
\choices{A. 多点接入、载波监听、碰撞检测}{B. 路由选择、流量控制、拥塞控制}
{C. 采样、量化、编码}{D. 封装、透明传输、差错检测}

\q 10 Mb/s 以太网的争用期是（\quad）
\choicesline{A. 5.12 $\mu$s}{B. 51.2 $\mu$s}{C. 512 $\mu$s}{D. 64 $\mu$s}

\q 10 Mb/s 以太网的最短帧长是（\quad）
\choicesline{A. 32 B}{B. 64 B}{C. 512 B}{D. 1518 B}

\q 截断二进制指数退避算法中，第 $i$ 次冲突后参数 $k$ 的取值为（\quad）
\choicesline{A. $k=i$}{B. $k=\min(i,10)$}{C. $k=\max(i,10)$}{D. $k=10$}

\q 以太网 MAC 地址的长度是（\quad）
\choicesline{A. 8 bit}{B. 16 bit}{C. 32 bit}{D. 48 bit}

\q 交换机自学习建立 MAC 地址表时，依据的是帧的（\quad）
\choicesline{A. 源 MAC 地址}{B. 目的 MAC 地址}{C. 类型字段}{D. FCS 字段}

\q IEEE 802.11 无线局域网的 MAC 层使用的协议是（\quad）
\choicesline{A. CSMA/CD}{B. CSMA/CA}{C. PPP}{D. HDLC}

\q 100BASE-T 中的“100”和“T”分别表示（\quad）
\choices{A. 100 Mb/s 和双绞线}{B. 100 m 和同轴电缆}
{C. 100 Mb/s 和光纤}{D. 100 个站点和双绞线}

\q 广域网中广泛使用的点对点数据链路层协议是（\quad）
\choicesline{A. PPP}{B. HTTP}{C. UDP}{D. CSMA/CA}

\qtypename{二、多项选择题}{3 题，每题 4 分，共 12 分；每题有两个或两个以上正确项}
\q 以太网 MAC 帧中包括的字段有（\quad）
\choicesline{A. 目的 MAC 地址}{B. 源 MAC 地址}{C. 类型字段}{D. 路由表}

\q 关于 CSMA/CD，下列说法正确的有（\quad）
\choices{A. 发送时边发送边检测冲突}{B. 检测到冲突后立即停止发送}
{C. 冲突后按退避算法等待并重发}{D. 全双工方式下必须使用 CSMA/CD}

\q 关于无线局域网，下列说法正确的有（\quad）
\choices{A. 遵循 IEEE 802.11 标准}{B. MAC 层使用 CSMA/CA 协议}
{C. 可借助 RTS/CTS 机制缓解隐藏站问题}{D. 与有线以太网一样使用 CSMA/CD}

\qtypename{三、判断题}{6 题，每题 2 分，共 12 分；正确的画“√”，错误的画“×”并写出改正后的表述}
\q 10 Mb/s 以太网的争用期是 512 $\mu$s。\quad（\quad）\quad 改正：\fixblank

\q 以太网最短帧长 64 B，长度小于 64 B 的帧均是冲突而中止的无效帧。\quad（\quad）\quad 改正：\fixblank

\q 交换机自学习 MAC 地址表时依据的是目的 MAC 地址。\quad（\quad）\quad 改正：\fixblank

\q VLAN 可把一个物理局域网划为多个广播域，提高安全性。\quad（\quad）\quad 改正：\fixblank

\q 全双工以太网仍然必须使用 CSMA/CD 协议。\quad（\quad）\quad 改正：\fixblank

\q 802.11 无线局域网使用 CSMA/CA，是因为无线环境中难以像有线以太网那样边发送边检测冲突。\quad（\quad）\quad 改正：\fixblank

\qtypename{四、计算题}{2 题，每题 5 分，共 10 分；要求写出公式和计算步骤}
\q 某 100 Mb/s 以太网采用 CSMA/CD 协议，信号在电缆中的传播速率为 $2\times10^{8}$ m/s，网络最大跨度为 2 km。求：（1）最大单程传播时延；（2）最小争用期（取往返传播时延）；（3）为在此争用期内检测到冲突，最短帧长应为多少比特、多少字节（写出计算过程）。

\q 某 10 Mb/s 以太网中，某站点连续发生 3 次冲突。求：（1）第 3 次冲突后参数 $k$ 的值；（2）随机数 $r$ 的取值范围；（3）可能的退避时间范围（用 $\mu$s 表示，取争用期 51.2 $\mu$s）。

\qtypename{五、简答题}{2 题，每题 8 分，共 16 分}
\q 简述 CSMA/CD 协议的工作过程，并说明争用期与最短帧长之间的关系。\anslines{2.4cm}

\q 简述以太网交换机自学习的工作原理（包括地址学习和帧的转发／过滤），并说明 VLAN 的作用。\anslines{2.4cm}

\newpage
\qtypename{六、综合题}{1 题，共 20 分}
\q 某农业园区建设局域网：办公楼、温室、仓库分别设置以太网交换机，通过光纤连接到核心交换机；温室内部署无线 AP 采集环境数据。网络按业务划分为办公、温室、视频监控三个 VLAN。请回答：（1）说明交换机自学习建立 MAC 地址表的过程，以及收到目的地址未知的帧时如何处理；（2）说明 VLAN 的作用，并指出不同 VLAN 之间通信需要什么设备；（3）温室无线终端接入采用什么协议？为什么不能像有线以太网那样使用 CSMA/CD？（4）若核心交换机到温室交换机的光纤长 10 km，链路速率 100 Mb/s，帧长 1000 B，求发送时延和传播时延（传播速率取 $2\times10^{8}$ m/s），并说明高速以太网为什么采用全双工方式；（5）从覆盖范围、核心技术、典型协议三方面比较局域网与广域网，说明园区网为什么选择以太网技术。\anslines{5cm}

\end{document}
